\section{Metodologia e plano de trabalho}
    A metodologia será organizada em três fases distintas, com detalhamento das atividades em cada uma:

    \subsection{Fase 1 - Fundamentação e Configuração do Ambiente (Semestre 1)}
        \begin{itemize}
            \item Levantamento de literatura em teste automatizado: pesquisa de artigos e livros sobre geração de testes em Java (EvoSuite [3], Randoop [10], PEX [4]). Análise de abordagens baseadas em algoritmos evolutivos, fuzzing e análise simbólica.
            \item Revisão sobre LLMs para geração de código e testes: leitura de trabalhos recentes (Chen et al., 2021 [1]; Chen et al., 2023 [2]; Moradi Dakhel et al., 2023 [8]; Gu et al., 2024 [5]; Steenhoek et al., 2024 [13]), identificando métodos, resultados e lacunas.
            \item Consolidação de métricas de cobertura: estudo de ferramentas JaCoCo [6], Cobertura [15] e EclEmma [14] para extração de métricas de cobertura de instruções, ramos, definições-uso e usos.
            \item Escolha de LLM(s): avaliação de modelos proprietários (OpenAI GPT-4, via API) e open-source (Code Llama, GPT-NeoX), observando custo, licenciamento e infraestrutura necessária.
            \item Configuração do ambiente de experimento: obtenção de chave API ou instalação local de modelo, instalação de bibliotecas Python (OpenAI SDK [9], Transformers [12], PyTorch [11]) e ferramentas Java (JUnit 5 [7], JaCoCo, Maven/Gradle).
            \item Criação de repositório versionado (Git): estruturação inicial com scripts de geração e diretórios para testes.
            \item Implementação de script-protótipo: script (Python ou Shell) que receba arquivo Java de um método, forme prompt e envie ao LLM, armazenando o teste JUnit gerado.
            \item Validação inicial de geração: testes com métodos simples (cálculo de fatorial, busca binária) para verificar compilação e execução.
            \item Mini-relatório de referência: documentação das descobertas teóricas, definindo estado da arte e justificando a proposta.
        \end{itemize}
    \subsection{Fase 2 - Seleção de Métodos-Alvo, Geração e Execução de Testes (Semestre 2)}
        \begin{itemize}
            \item Seleção de 30 métodos Java: escolha em projetos open-source de métodos com diversos padrões de controle (condicionais aninhados, laços, recursão, streams, exceções).
            \item Análise manual de cobertura esperada: mapeamento de todos os caminhos lógicos e pontos de definições-uso de variáveis para cada método, incluindo cálculo da complexidade ciclomática.
            \item Documentação em planilha: registro de nomes de métodos, complexidade, número de caminhos e variáveis críticas para def-use e dados auxiliares.
            \item Construção de prompts de geração: elaboração de instruções em linguagem natural detalhando o método-fonte e solicitando a criação de testes JUnit que cubram todos os caminhos ou definições-uso.
            \item Execução automatizada em lote: pipeline que percorra a lista de métodos, envie prompts ao LLM e armazene arquivos \texttt{.java} com classes de teste JUnit correspondentes.
            \item Primeira medição de cobertura: integração do plugin JaCoCo no build, execução de todos os testes e geração de relatórios (HTML, XML, CSV) contendo métricas de cobertura de instruções e ramos.
            \item Refinamento iterativo de prompts: ajustes nas instruções para lidar com cenários específicos (cobertura de exceções, valores de entrada extremados, variáveis críticas).
            \item Configuração de projeto Java: estrutura Maven/Gradle que inclua código-fonte original e testes gerados, possibilitando compilação e execução unificada via comando.
        \end{itemize}
    \subsection{Fase 3 - Avaliação, Documentação e Disseminação (Semestre 3)}
        \begin{itemize}
            \item Análise de relatórios de cobertura: comparativo entre cobertura obtida pelo LLM e cobertura teórica (documentada na fase 2); identificação de lacunas (caminhos não cobertos, usos faltantes).
            \item Comparação com ferramentas tradicionais: execução do EvoSuite [3] para cada método, geração de testes automáticos e medição de cobertura via JaCoCo; compilação de tabela comparativa entre LLM e EvoSuite.
            \item Refinamento de prompts: para métodos com cobertura insuficiente, ajustes nos prompts (inclusão de exemplos, templates restritos) e nova execução do pipeline.
            \item Avaliação qualitativa dos testes gerados: análise de legibilidade, aderência a boas práticas de JUnit (uso correto de anotações, asserts descritivos) e presença de ``test smells''; registro de conclusões em relatório.
            \item Consolidação de dados: organização de resultados quantitativos (planilhas, gráficos comparativos de cobertura) e elaboração de considerações iniciais.
            \item Redação do relatório técnico final: capítulos: Introdução, Fundamentação Teórica, Metodologia, Resultados e Discussão, Conclusão e Trabalhos Futuros.
        \end{itemize}